\section*{Capítulo \ref{ch:g5:areas} --- Perímetro e área}

\begin{solution}{exo:g5:areas:1}
\omterm{def:g4:measure:area}{Área}: $8 \times 5 = 40$ cm$^2$. \omterm{def:g4:measure:perimeter}{Perímetro}:
$2 \times (8 + 5) = 26$ cm.
\end{solution}

\begin{solution}{exo:g5:areas:2}
$9 \times 9 = 81$ cm$^2$; \quad $12 \times 7 = 84$ m$^2$; \quad
$4.5 \times 6 = 27$ cm$^2$.
\end{solution}

\begin{solution}{exo:g5:areas:3}
Largura: $63 \div 9 = 7$ cm. \omterm{def:g4:measure:perimeter}{Perímetro}: $2 \times (9 + 7) = 32$ cm.
\end{solution}

\begin{solution}{exo:g5:areas:4}
Por exemplo $6 \times 4$ (\omterm{def:g4:measure:perimeter}{perímetro} $20$) e $8 \times 3$ (\omterm{def:g4:measure:perimeter}{perímetro}
$22$): mesma \omterm{def:g4:measure:area}{área} $24$, \omterm{def:g4:measure:perimeter}{perímetros} diferentes.
\end{solution}

\begin{solution}{exo:g5:areas:5}
$3$ m$^2 = 30\,000$ cm$^2$; \quad $50\,000$ cm$^2 = 5$ m$^2$.
\end{solution}

\begin{solution}{exo:g5:areas:6}
Barra: $8 \times 2 = 16$ cm$^2$; haste: $2 \times 5 = 10$ cm$^2$;
total $26$ cm$^2$.
\end{solution}

\begin{solution}{exo:g5:areas:7}
$30 \times 20 - 24 \times 14 = 600 - 336 = 264$ cm$^2$ de madeira.
\end{solution}

\begin{solution}{exo:g5:areas:8}
\omterm{def:g3:division:half}{Metade} do retângulo $8 \times 6$: $48 \div 2 = 24$ cm$^2$.
\end{solution}

\begin{solution}{exo:g5:areas:9}
Semente: \omterm{def:g4:measure:area}{área} $7 \times 4 = 28$ m$^2$, custo $28 \times 2 = 56$.
Cerca: \omterm{def:g4:measure:perimeter}{perímetro} $2 \times (7 + 4) = 22$ m, custo
$22 \times 3 = 66$.
\end{solution}

\begin{solution}{exo:g5:areas:10}
Dois ladrilhos de $50$ cm fazem um metro: ao longo dos $12$ m de
comprimento, $24$ ladrilhos; ao longo dos $2$ m de largura, $4$
ladrilhos. Total: $24 \times 4 = 96$ ladrilhos.
\end{solution}

\begin{solution}{exo:g5:areas:11}
Antes: \omterm{def:g4:measure:perimeter}{perímetro} $2 \times (4 + 3) = 14$ cm, \omterm{def:g4:measure:area}{área} $12$ cm$^2$.
Depois (lados $8$ e $6$): \omterm{def:g4:measure:perimeter}{perímetro} $28$ cm, \omterm{def:g4:measure:area}{área} $48$ cm$^2$. O
\omterm{def:g4:measure:perimeter}{perímetro} \emph{dobrou}; a \omterm{def:g4:measure:area}{área} foi multiplicada por \emph{quatro}
($2 \times 2$): comprimentos e \omterm{def:g4:measure:area}{áreas} não mudam de escala do mesmo
jeito.
\end{solution}
